\documentclass{article}

\usepackage{amsmath,amssymb}
\usepackage{graphicx}

\title{ICRA2027 - WAMPlan}
\author{Aleksandr Panov}
\date{September 2026}

\begin{document}

\maketitle

\section{Introduction}

\section{Related Work}

\section{Preliminaries}

\section{Consensus Planning}
\label{sec:method}

Large generative action models can represent multiple plausible continuations
from the same observation. This stochasticity, however, creates a practical
decision problem at test time: a single sampled action chunk may be locally
implausible even when other samples agree on a competent continuation. We
introduce \emph{consensus planning}, a training-free test-time procedure that
selects a representative action chunk from a small set of samples. The method
operates entirely in action space, requires neither additional supervision nor
access to a value function, and can therefore be applied to a pretrained
generative policy without modifying its parameters.

\subsection{Action-Chunk Generation}

At control query $q$, the policy receives the current observation history and a
language instruction, denoted jointly by $x_q$. It produces a horizon-$T$
action chunk
\begin{equation}
    \mathbf{a}^{(i)} = \left(a^{(i)}_1,\ldots,a^{(i)}_T\right),
    \qquad \mathbf{a}^{(i)} \sim p_\theta(\mathbf{a}\mid x_q; s_i),
    \label{eq:candidates}
\end{equation}
where $s_i$ is the sampling seed and $i\in\{1,\ldots,K\}$ indexes a candidate.
In our implementation, $K=3$ and the candidate seeds are fixed to $\{0,1,2\}$
for reproducible test-time selection. This seed controls the stochastic
generation of an action chunk; it is distinct from the benchmark episode seed,
which specifies the scene instance.

Each action $a_t$ contains a translational command $p_t\in\mathbb{R}^3$, a
continuous 6D rotation representation $r_t\in\mathbb{R}^6$, and a scalar
gripper command $g_t$. Although a chunk contains $T$ proposed actions, the
controller executes only its first $H$ actions before requesting a new chunk.
Consequently, consensus is evaluated on this executable prefix rather than on
long-horizon predictions that will be discarded by closed-loop replanning.

\subsection{Control-Aware Distance Between Chunks}

We compare candidates using a distance aligned with the robot control
parameterization. Let $R(r)\in SO(3)$ denote the rotation matrix obtained from
the 6D representation by orthonormalizing its two three-dimensional vectors.
For two individual actions $a$ and $b$, we define
\begin{equation}
\begin{split}
    d_{\mathrm{act}}(a,b) ={}& \lambda_p\frac{\lVert p_a-p_b\rVert_2}{\sqrt{3}}
    + \lambda_r\frac{\arccos\!\left(
        \operatorname{clip}\!\left(
        \frac{\operatorname{tr}(R(r_a)R(r_b)^\top)-1}{2},-1,1
        \right)\right)}{\pi} \\
    &+ \lambda_g\,\mathbf{1}\!\left[\operatorname{sign}(g_a)
        \neq \operatorname{sign}(g_b)\right],
    \label{eq:action-distance}
\end{split}
\end{equation}
where the first term measures normalized translational disagreement, the second
is the normalized geodesic distance on $SO(3)$, and the third penalizes an
incompatible gripper mode. We use $\lambda_p=1$, $\lambda_r=0.5$, and
$\lambda_g=0.25$. These terms prevent a high-dimensional Euclidean distance
from conflating translation, orientation, and discrete gripper decisions.

The distance between two candidate chunks is a temporally discounted average
over their executable prefixes:
\begin{equation}
    D_H\!\left(\mathbf{a}^{(i)},\mathbf{a}^{(j)}\right)
    =
    \frac{\sum_{t=1}^{H}\gamma^{t-1}
    d_{\mathrm{act}}\!\left(a^{(i)}_t,a^{(j)}_t\right)}
    {\sum_{t=1}^{H}\gamma^{t-1}},
    \label{eq:chunk-distance}
\end{equation}
with $\gamma=0.9$. The discount assigns greater importance to actions that
will be executed sooner and whose errors cannot be corrected before the next
query.

\subsection{Medoid Selection and Closed-Loop Execution}

Consensus planning selects the candidate with minimum mean disagreement to the
remaining candidates:
\begin{equation}
    i^* = \arg\min_{i\in\{1,\ldots,K\}}
    \frac{1}{K-1}\sum_{\substack{j=1 \\ j\neq i}}^K
    D_H\!\left(\mathbf{a}^{(i)},\mathbf{a}^{(j)}\right).
    \label{eq:medoid}
\end{equation}
This is a medoid rather than an action-space average: the selected chunk is an
actual sample from the generative policy. The distinction is important for
robot actions, since averaging two valid but qualitatively different motions can
produce an invalid intermediate command, especially near gripper transitions or
rotational discontinuities.

The controller executes the prefix
$\left(a^{(i^*)}_1,\ldots,a^{(i^*)}_H\right)$, observes the resulting state,
and repeats Eqs.~\eqref{eq:candidates}--\eqref{eq:medoid}. Thus, consensus is
not a commitment to a full predicted trajectory: it is a receding-horizon
selection rule that uses sample agreement as a local reliability signal. When
the candidates are close, the procedure retains a representative prediction;
when one candidate is an outlier, its larger average distance makes it unlikely
to be executed. The method does not assume that agreement proves task success,
and its empirical utility must therefore be established by the evaluations in
the subsequent sections.

\end{document}
